\section{Primitives, State Space, and Regularity Conditions}

This phase freezes the benchmark model before any static or dynamic optimization is derived. The purpose is to eliminate notation drift and to state all regularity requirements used later in Phase 2 (static controls), Phase 3 (Bellman fixed point), and Phase 4 (measurable choice regions and transition matrix objects).

\subsection{Benchmark primitives}

Time is discrete, with periods $t=0,1,2,\dots$.

The set of active organizational types is
\[
\mathcal T=\{A,B,C\},
\]
where $A$ denotes integrated firms, $B$ denotes research-specialized firms, and $C$ denotes manufacturing-specialized firms.

Medicine categories are
\[
k\in\{G,O\},
\]
where $G$ means generic and $O$ means original drugs.

The policy regime is binary:
\[
M\in\{0,1\},
\]
with $M=0$ interpreted as pre-MAH and $M=1$ as post-MAH.

Each incumbent firm carries a one-dimensional capability state
\[
z\in Z\subset\mathbb R_+.
\]
The benchmark intentionally uses one idiosyncratic state dimension to keep the model identifiable and computationally implementable at this stage.

Prices are
\[
(w,p_m),
\]
where $w$ is wage (or a composite factor price) and $p_m$ is the qualified contract-manufacturing service price.

For each chosen organizational type $j\in\{A,B,C\}$, next-period capability evolves through the kernel
\[
Q_j(\cdot\mid z,M).
\]

Entrants draw initial capability from probability measure $G$ on $Z$.

\paragraph{MAH channels in reduced form.}
The benchmark allows MAH to enter only through:
\begin{align}
\tau_M(M) &\quad \text{commercialization friction},\\
\kappa_{sj}(M) &\quad \text{switching cost from current type } s \text{ to chosen type } j.
\end{align}

For type-$B$ firms, commercialization of one original-drug idea is summarized by bounded reduced-form success rate
\[
\lambda_B(z,p_m,M)\in[0,1].
\]
Hence the one-idea value object is fixed as
\[
H_B(z;p_m,M)=\lambda_B(z,p_m,M)\bigl(R_B(z)-p_m\bigr)-\tau_M(M).
\]

\subsection{Timing and value mapping}

At the beginning of period $t$, a firm with current type $s$ and state $z$ chooses organizational form $j\in\{A,B,C,\text{exit}\}$. If $j\neq s$, it pays switching cost $\kappa_{sj}(M)$. Then it receives optimized flow payoff of chosen type $j$, denoted $\Pi_j(z;w,p_m,M)$, and transitions according to $Q_j(\cdot\mid z,M)$ if it does not exit.

Therefore choice-specific value is
\[
W_{sj}(z;M)=\Pi_j(z;w,p_m,M)-\kappa_{sj}(M)+\beta\int V_j(z';M)Q_j(dz'\mid z,M),
\]
and
\[
W_{s,\text{exit}}(z;M)=0.
\]

The Bellman equations in Phase 3 will be written as
\[
V_s(z;M)=\max\{W_{sA}(z;M),W_{sB}(z;M),W_{sC}(z;M),0\},\qquad s\in\{A,B,C\}.
\]
This representation is fixed in Phase 1 and should not be reparameterized in later sections.

\subsection{State space, sigma algebra, and function spaces}

Let $\mathcal B(Z)$ be the Borel sigma algebra on $Z$. For each $j$, the map
\[
z\mapsto Q_j(B\mid z,M)
\]
is measurable for every $B\in\mathcal B(Z)$, and $Q_j(\cdot\mid z,M)$ is a probability measure on $(Z,\mathcal B(Z))$.

Define
\[
\mathcal C_b(Z)=\{f:Z\to\mathbb R\mid f \text{ continuous and bounded}\},
\qquad
\|f\|_\infty=\sup_{z\in Z}|f(z)|.
\]

Value-function domain is the product Banach space
\[
\mathcal V=\mathcal C_b(Z)^3,
\]
equipped with norm
\[
\|(V_A,V_B,V_C)\|_{\infty,3}=\max\{\|V_A\|_\infty,\|V_B\|_\infty,\|V_C\|_\infty\}.
\]

This is the exact space used for contraction mapping in Phase 3.

\subsection{Regularity assumptions}

\begin{assumption}[A1: Compact state support]
$Z$ is nonempty, compact, and metrizable.
\end{assumption}

\begin{assumption}[A2: Discounting]
$\beta\in(0,1)$.
\end{assumption}

\begin{assumption}[A3: Flow payoff regularity]
For each $j\in\{A,B,C\}$ and fixed $(w,p_m,M)$, the optimized flow payoff $\Pi_j(z;w,p_m,M)$ is bounded and continuous in $z$.
\end{assumption}

\begin{assumption}[A4: Weak continuity of kernels]
For each $j\in\{A,B,C\}$, fixed $M$, and each $f\in\mathcal C_b(Z)$,
\[
z\mapsto \int f(z')Q_j(dz'\mid z,M)
\]
is continuous on $Z$.
\end{assumption}

\begin{assumption}[A5: Convex static technologies]
Type-$A$ and type-$B$ innovation cost functions are convex in effort with elasticity strictly above one. Manufacturing cost $c(m,z;w)$ is continuous in $(m,z,w)$, strictly convex in $m$, and coercive in $m$.
\end{assumption}

\begin{assumption}[A6: Bounded MAH channels]
$\tau_M(M)$ and $\kappa_{sj}(M)$ are finite for all relevant indices; $\lambda_B(z,p_m,M)$ is measurable in $z$ and bounded in $[0,1]$.
\end{assumption}

\subsection{Why these assumptions are sufficient for later phases}

The following implications are used explicitly in Phase 2 and Phase 3.

\paragraph{Implication 1: static well-definedness.}
Under A3, A5, and A6, each static optimization problem has a finite value; strict convexity/coercivity ensures existence of maximizers, and strict convexity in control implies uniqueness when interior first-order conditions apply.

\paragraph{Implication 2: Bellman operator preserves continuity.}
Under A3 and A4, if $V_j\in\mathcal C_b(Z)$ then
\[
z\mapsto \Pi_j(z;w,p_m,M)+\beta\int V_j(z')Q_j(dz'\mid z,M)
\]
is continuous and bounded; therefore each $W_{sj}(\cdot;M)$ is in $\mathcal C_b(Z)$.

\paragraph{Implication 3: contraction structure.}
Under A2,
\[
\|\mathcal T(V)-\mathcal T(\tilde V)\|_{\infty,3}
\le
\beta\|V-\tilde V\|_{\infty,3},
\]
where $\mathcal T$ is the Bellman operator defined in Phase 3. Because $\beta<1$, Banach fixed-point theorem applies.

\paragraph{Implication 4: measurable choice regions.}
Because $W_{sj}(\cdot;M)$ are continuous and thus measurable, argmax correspondence over a finite action set is measurable; this is the key object needed to define choice regions and transition probabilities in Phase 4.

\subsection{Benchmark boundary statement}

The benchmark is frozen by four non-negotiable constraints:
\begin{enumerate}
    \item Active types remain exactly $A$, $B$, and $C$.
    \item State remains one-dimensional in firm capability $z$, plus policy state $M\in\{0,1\}$.
    \item Type-$B$ commercialization remains reduced form via $\lambda_B(z,p_m,M)$; no explicit top-level implementation-intensity control is introduced.
    \item MAH enters benchmark primitives through $\tau_M(M)$ and $\kappa_{sj}(M)$, without adding government-optimal-policy or DSGE layers.
\end{enumerate}

These objects and assumptions are the fixed base for the next phase: static operating payoffs and optimal controls.
